\documentclass[12pt]{article}
\usepackage[utf8]{inputenc}
\usepackage[T1]{fontenc}
\usepackage{hyperref}
\title{Yami Language Resources for K-12 Curriculum Development}
\author{Compiled by Claude Sonnet 5 (claude-sonnet-5) for Andr\'as Kornai}
\date{}
\begin{document}
\maketitle
\section{Introduction}
Yami, also called Tao (ISO 639-3 \texttt{tao}), is an Austronesian language of Orchid Island, Taiwan. Sources give roughly 4,000 speakers, and the older figure is 3,800. UNESCO is reported to classify it as Vulnerable.
\section{Primary Educational Resources}
\begin{itemize}
\item Tao (Yami) Language e-Learning, a virtual classroom begun in 2005 at Providence University by Victoria Rau and Meng-Chien Yang. According to the search summary, since September 2025 it has been moving to National Chung Cheng University's ALARAD Center. The lessons follow the standards of Taiwan's Ministry of Education: \url{https://taotech.url.tw/taoelearning/en/pages/home.html}
\end{itemize}
\section{Linguistic and Reference Materials}
\begin{itemize}
\item \emph{Yami Texts with Reference Grammar and Dictionary} (Academia Sinica, 2006), as listed in search summaries; it is available online through the Providence site. See: \url{https://pinyin.info/news/2006/tao-yami-language-materials/}
\item Published dictionaries reported in search summaries: \emph{Yami Learning Dictionary} (2008), \emph{Yami (Tao) Dictionary} (NTU Press, 2012, Chinese/English) and \emph{Yami Bible Dictionary} (2018). The web pages were not individually opened.
\item Yami Documentation archive: \url{https://eldp.access.preservica.com/dk0110/}
\item Endangered Languages Project entry: \url{https://www.endangeredlanguages.com/lang/10747?hl=zh-cn}
\end{itemize}
\section{Religious and Community Resources}
The \emph{Yami Bible Dictionary} (2018) is listed above. No other religious or community resources were examined.
\section{Digital and Audio Resources}
The e-Learning site above and the ELDP archive of documentation recordings: \url{https://eldp.access.preservica.com/dk0110/}
\section{Curriculum Development Context}
Schooling is reportedly mainly in Chinese. Search summaries say Yami has been an elective in elementary school since 1998. The e-Learning course and the dictionaries are the best base for a native-language curriculum.
\end{document}